% ============================================================
% CHAPTER 5 - IMPLEMENTATION AND TESTING
% ============================================================
\chapter{Implementation and Testing}
\label{ch:implementation}

\section{Introduction}

This chapter describes how the design of Chapter~\ref{ch:analysis-design} was
realised. Section~\ref{sec:impl-env} states the development and execution
environments. Sections~\ref{sec:impl-training} and \ref{sec:impl-service}
describe the training pipeline and the inference service.
Section~\ref{sec:impl-challenges} documents the implementation problems that
materially affected the outcome, several of which changed the design.
Section~\ref{sec:impl-frontend} covers the web client and
Section~\ref{sec:impl-deploy} the deployment. Section~\ref{sec:testing} sets out
the testing strategy and results.

\section{Development and Execution Environments}
\label{sec:impl-env}

Model development was carried out in a Kaggle notebook with an NVIDIA Tesla T4
GPU, chosen because the wav2vec2 alignment pass over 20{,}000 recordings is
GPU-bound and impractical on CPU. The application was developed locally on
Windows 11 and deployed as a Docker image. Table~\ref{tab:environments}
summarises the three environments and the reason each exists.

\begin{table}[H]
\centering
\caption[Execution environments]{Execution environments.}
\label{tab:environments}
\small
\begin{tabularx}{\textwidth}{@{}p{2.6cm}p{3.6cm}X@{}}
\toprule
\textbf{Environment} & \textbf{Hardware} & \textbf{Purpose} \\
\midrule
Kaggle notebook & Tesla T4 GPU, 16\,GB & Dataset construction, training, architecture comparison \\
Local workstation & Windows 11, CPU & Application development, integration testing \\
Container / Space & 2 vCPU, 16\,GB, CPU only & Production deployment of the inference service and UI \\
\bottomrule
\end{tabularx}
\end{table}

\section{Training Pipeline Implementation}
\label{sec:impl-training}

The training pipeline is implemented as a sequence of notebook stages
corresponding directly to Sections~\ref{sec:g2p}--\ref{sec:split}.

\subsection{Grapheme-to-phoneme conversion}

G2P is a subprocess call to the \texttt{espeak-ng} binary with diacritic
stripping, followed by a vocabulary filter
(Listing~\ref{lst:g2p}). The three stripped symbols are given as Unicode escapes
so that the source file remains ASCII.

\begin{lstlisting}[language=Python, caption={Grapheme-to-phoneme conversion via the espeak-ng binary.}, label={lst:g2p}]
def word_to_phonemes(word):
    """IPA phonemes for one word, via the espeak-ng CLI (not the
    phonemizer wrapper, whose library-copy step fails when /tmp is
    mounted noexec)."""
    out = subprocess.run(
        ["espeak-ng", "-q", "--ipa", "-v", "en-us", word.strip()],
        capture_output=True, text=True,
    ).stdout.strip()
    # Strip primary stress (U+02C8), secondary stress (U+02CC) and the
    # length mark (U+02D0): not members of the aligner's phone vocabulary.
    for mark in ("\u02c8", "\u02cc", "\u02d0"):
        out = out.replace(mark, "")
    return [ch for ch in out if ch.strip()]
\end{lstlisting}

\subsection{Forced alignment}

Alignment loads the recording, trims it once, computes the CTC emission matrix
and solves Equation~\eqref{eq:ctc-align}
(Listing~\ref{lst:align}). The frame-to-second ratio is derived from the
\emph{trimmed} sample count so that the returned timestamps share a timebase
with everything downstream.

\begin{lstlisting}[language=Python, caption={CTC forced alignment of a waveform to a target phoneme sequence.}, label={lst:align}]
def align_phonemes(audio_path, transcript):
    waveform, _ = librosa.load(audio_path, sr=SAMPLE_RATE)
    waveform, _ = librosa.effects.trim(waveform, top_db=TRIM_TOP_DB)
    if len(waveform) < SAMPLE_RATE * 0.05:
        raise ValueError("audio too short after trimming (<50ms)")
    waveform_t = torch.tensor(waveform).unsqueeze(0)

    # Keep only phonemes the aligner actually knows about.
    tokens, kept = [], []
    for p in transcript_to_phonemes(transcript):
        if p in _VOCAB:
            tokens.append(_VOCAB[p]); kept.append(p)
    if not tokens:
        raise ValueError("no valid phonemes for '%s'" % transcript)

    emission = _emission(waveform_t)                  # log P(k | t)
    targets = torch.tensor([tokens], dtype=torch.int32, device=device)
    aligned, scores = torchaudio.functional.forced_align(
        emission, targets, blank=BLANK_ID)
    aligned, scores = aligned[0], scores[0].exp()
    spans = torchaudio.functional.merge_tokens(aligned, scores)

    # Frame index -> seconds, on the TRIMMED timebase.
    ratio = waveform_t.size(1) / emission.size(1) / SAMPLE_RATE
    return [(p, float(s.start * ratio), float(s.end * ratio), float(s.score))
            for s, p in zip(spans, kept[: len(spans)])]
\end{lstlisting}

\subsection{Dataset construction and control-baseline labelling}

Dataset construction is a two-pass procedure, and the two-pass structure is
required by the labelling rule: the threshold $\tau$ of
Equation~\eqref{eq:method2} is a percentile of the dysarthric deviation
distribution and therefore cannot be known until every deviation has been
computed. The first pass aligns each pair, crops both patches, computes $D_i$
and records the provenance of every instance; the second pass computes $\tau$
and applies the labelling rule
(Listing~\ref{lst:labelling}).

\begin{lstlisting}[language=Python, caption={Control-baseline labelling (Method 2). The cutoff is computed over the dysarthric distances only.}, label={lst:labelling}]
# Pass 1 (per pair): collect BOTH the dysarthric patch (label decided
# later) and the matched control patch (correct by definition).
for i, (phon, s, e, sc) in enumerate(dys_align):
    if i >= len(ctrl_align):
        continue                       # no matched control phoneme
    _, cs, ce, csc = ctrl_align[i]
    imp_mel  = crop_phoneme_mel(dys_path,  s,  e)
    ctrl_mel = crop_phoneme_mel(ctrl_path, cs, ce)

    mels.append(imp_mel);  dists.append(phoneme_distance(imp_mel, ctrl_mel))
    is_control.append(0)               # dysarthric instance
    mels.append(ctrl_mel); dists.append(0.0)
    is_control.append(1)               # control anchor, always correct

# Pass 2: threshold estimated over the DYSARTHRIC distances ONLY, so a
# control instance can never be dragged into the error class.
dys_mask = ctrl_arr == 0
cutoff = np.percentile(dist_arr[dys_mask], 100 - ERROR_PERCENTILE)

y = np.zeros(len(X), dtype=int)
y[dys_mask & (dist_arr >= cutoff)] = 1
\end{lstlisting}

The contrast with Method~1 is a single line. Method~1 computed the percentile as
\pth{np.percentile(dist_arr, 100 - ERROR_PERCENTILE)} over the whole array,
without the \pth{dys_mask}; that omission is the entire defect described in
Section~\ref{sec:method1}.

\subsection{Split, balancing and training}

The ordering constraint of Section~\ref{sec:split} is enforced explicitly:
\texttt{train\_test\_split} with stratification runs first, and
\texttt{sklearn.utils.resample} is applied afterwards to the training partition
only. The model is constructed with the Keras functional API exactly as
specified in Table~\ref{tab:model-layers}, compiled with Adam and categorical
cross-entropy, and fitted with early stopping on validation loss with a patience
of five epochs and best-weight restoration.

\section{Inference Service Implementation}
\label{sec:impl-service}

\subsection{Structure}

The service is a FastAPI application of six modules:
\texttt{main} (routing and lifecycle), \texttt{pipeline} (models and
inference), \texttt{audio} (decoding, trimming, quality assessment,
downsampling), \texttt{phonemes} (G2P), \texttt{config} (typed settings) and
\texttt{schemas} (response models).

\subsection{Model lifecycle}

The three heavy models total approximately 1.3\,GB and take tens of seconds to
load. They are therefore loaded once and shared. A process-wide singleton holds
the pipeline; loading is lazy and guarded by a double-checked lock so that
concurrent first requests cannot trigger duplicate loads; and a FastAPI lifespan
hook warms the models at start-up so the first user request is not the one that
pays the cost. Warm-up failure is caught and logged rather than propagated,
leaving the service running in a degraded state that the health endpoint
reports (NFR-09, NFR-13).

\subsection{Normalisation-constant resolution}
\label{sec:impl-normalisation}

The correctness requirement NFR-03 reduces in practice to a single question:
does the service apply the same $\mu$ and $\sigma$ that training applied? The
resolution order implemented is environment variables, then the dataset archive
if present, then the constants compiled into \texttt{config.py}
(Listing~\ref{lst:norm}). The fallback to per-clip normalisation exists only for
the case where all three sources are absent, and it emits a warning, because a
service that silently degrades its own accuracy is worse than one that fails
loudly.

\begin{lstlisting}[language=Python, caption={Resolution of the global mel normalisation constants.}, label={lst:norm}]
# Decoded once from the shipped phoneme_dataset.npz. These two scalars are
# as much a part of the trained model as its weights: the log-mel features
# are in dB, so the input scaling must match training exactly.
_KNOWN_MEL_MEAN = -2.41279196739197
_KNOWN_MEL_STD  = 10.222767829895

def load_norm_stats(self):
    """Precedence: env override -> dataset npz -> baked-in constants."""
    if self.mel_mean is not None and self.mel_std is not None:
        return self.mel_mean, self.mel_std
    if self.dataset_npz_path.exists():
        try:
            d = np.load(self.dataset_npz_path, allow_pickle=True)
            if "mel_mean" in d and "mel_std" in d:
                return float(d["mel_mean"]), float(d["mel_std"])
        except Exception:
            pass
    return _KNOWN_MEL_MEAN, _KNOWN_MEL_STD
\end{lstlisting}

\subsection{Quality gate}

The gate computes the RMS level in dBFS and the mean spectral flatness of the
trimmed waveform and rejects the request with a user-facing message if either
fails its bound (Listing~\ref{lst:quality}). Spectral flatness --- the ratio of
the geometric to the arithmetic mean of the power spectrum --- approaches unity
for white noise and is small for voiced speech, which makes it an inexpensive
and effective discriminator between a genuine attempt and room tone.

\begin{lstlisting}[language=Python, caption={Pre-inference audio quality gate.}, label={lst:quality}]
def assess_quality(wave, min_rms_dbfs, max_flatness):
    """Returns a human-readable rejection reason, or None if usable."""
    if len(wave) == 0:
        return ("No audio detected. Please record again and make sure "
                "your microphone is working.")

    rms  = float(np.sqrt(np.mean(np.square(wave)) + 1e-12))
    dbfs = 20 * np.log10(rms + 1e-12)
    if dbfs < min_rms_dbfs:                       # default -40 dBFS
        return ("Recording is too quiet. Please speak louder or move "
                "closer to the microphone.")

    flatness = float(np.mean(librosa.feature.spectral_flatness(y=wave)))
    if flatness > max_flatness:                   # default 0.6
        return ("This sounds like background noise rather than speech. "
                "Please try recording again in a quieter space.")
    return None
\end{lstlisting}

\subsection{Target-word verification}

Verification transcribes the recording with \texttt{whisper-tiny.en}, normalises
both the transcript and the target with Whisper's own text normaliser --- which
reconciles spelled-out numbers with digits, strips punctuation and folds case,
so that a target of ``seven'' matches a transcription of ``7'' --- and computes
the maximum \texttt{SequenceMatcher} ratio between the target and any
transcribed word. A match is declared if that ratio meets the threshold or if
either string contains the other. The check is intentionally permissive: its
purpose is to catch the wrong-word case, and a false rejection of a genuine but
heavily distorted attempt would be far more harmful than an occasional missed
detection. Any exception in this stage is caught and analysis proceeds.

\subsection{Analysis}

The analysis method is a direct transcription of the notebook inference cell:
align, crop one mel patch per phoneme, stack, normalise, and classify the whole
batch in a single \texttt{predict} call. Batching matters --- a five-phoneme word
issues one framework call rather than five --- and the per-phoneme threshold is
applied afterwards in pure Python, which is what makes the user-adjustable
threshold free of any additional inference cost.

\section{Implementation Challenges and Resolutions}
\label{sec:impl-challenges}

Six problems arose that materially affected the result. They are recorded here
because several of them are the kind of defect that produces a working system
with quietly wrong numbers.

\subsection{The phonemizer library failure}

\emph{Symptom.} \texttt{phonemizer} raised ``failed to map segment from shared
object'' on the training host.

\emph{Cause.} The package copies the \texttt{espeak} shared library into a
temporary directory before loading it; the training environment mounts the
temporary filesystem \texttt{noexec}, so the copied library cannot be executed.

\emph{Resolution.} The \texttt{espeak-ng} binary is invoked directly through a
subprocess, which runs the installed library in place. The same code path is
used in training and in the service, which additionally eliminates a source of
train--serve skew. A start-up self-check now converts the word ``well'' and
fails loudly if the result is empty or disjoint from the aligner's vocabulary,
so a broken toolchain is detected before any data is processed rather than after
hours of silently empty output.

\subsection{The double-trim timebase defect}

\emph{Symptom.} Phoneme bands drawn on the waveform were displaced relative to
the audible sounds.

\emph{Cause.} Alignment was computed on the trimmed waveform, while the
displayed waveform was loaded separately and untrimmed. The two timebases
differed by the duration of the leading silence.

\emph{Resolution.} Load-and-trim was factored into a single function whose
output is the sole waveform used by alignment, cropping and display. This
``trim once'' discipline is now a stated invariant of both the notebook and the
service.

\subsection{Test-set contamination through premature balancing}

\emph{Symptom.} Implausibly high held-out accuracy in an early iteration.

\emph{Cause.} Minority-class upsampling was applied before the train/test split,
so identical duplicated instances appeared on both sides of it. The model was
being evaluated partly on data it had memorised.

\emph{Resolution.} The split is performed first, with stratification, and
balancing is applied to the training partition only. The held-out partition
retains the natural 15\%/85\% distribution. Reported accuracy fell substantially
and became meaningful.

\subsection{Control speech flagged as erroneous}

\emph{Symptom.} Clear control recordings were returned as full of errors.

\emph{Cause.} Method~1's global percentile labelling
(Section~\ref{sec:method1}).

\emph{Resolution.} The control-baseline labelling of Method~2. This is the
principal methodological contribution of the dissertation and is treated in full
in Section~\ref{sec:method2}; it is listed here because it was discovered as an
implementation-stage empirical failure rather than as a theoretical objection.

\subsection{Recurrent dropout disabling the fused kernel}

\emph{Symptom.} Training time per epoch roughly an order of magnitude higher
than expected on the GPU.

\emph{Cause.} Specifying \texttt{recurrent\_dropout} on the LSTM layer
disqualifies it from the fused cuDNN implementation, forcing a slow generic
loop.

\emph{Resolution.} Recurrent dropout was removed and regularisation retained
through input dropout on the recurrent layer and dropout after each
convolutional block. Epoch time returned to approximately 50\,s and no loss of
generalisation was observed.

\subsection{Browser audio format mismatch}

\emph{Symptom.} Recordings made in the browser could not be decoded server-side.

\emph{Cause.} \texttt{MediaRecorder} emits WebM/Opus; \texttt{libsndfile} does
not decode it, and the \texttt{librosa} fallback path requires \texttt{ffmpeg}.

\emph{Resolution.} Conversion was moved to the client (Web Audio API decode,
\texttt{OfflineAudioContext} resample to 16\,kHz mono, hand-written PCM16 WAV
encoder), removing the native dependency from the server image and guaranteeing
that the server receives exactly the format training assumed.

\section{Web Client Implementation}
\label{sec:impl-frontend}

The client is a React 18 application built with Vite, deliberately free of a
component library or state-management framework: the interaction is a single
form and a single result view, and the entire state is six values held in the
root component.

Audio capture requests a mono stream with echo cancellation, noise suppression
and automatic gain control enabled. The last of these is a considered
compromise: gain control alters the level of the recording, which the model is
sensitive to through the global normalisation, but the alternative --- recordings
whose level varies with microphone distance and is frequently low enough to trip
the quality gate --- proved worse in practice.

The result view renders three coordinated presentations of the same result: the
radial clarity ring, the waveform with verdict-shaded phoneme spans drawn as an
inline SVG, and the phoneme list with per-phoneme probability meters. The
waveform is drawn from the peak-preserving decimation returned by the server
(at most 2000 points), which keeps the payload small while retaining the visual
envelope of the signal.

The application is installable as a progressive web application through a web
manifest, so a user may add it to a phone home screen. Microphone capture
requires a secure context, which is satisfied by \texttt{localhost} in
development and by HTTPS in deployment.

\section{Deployment}
\label{sec:impl-deploy}

Two deployment topologies are supported.

\textbf{Split containers (development).} \texttt{docker-compose} builds the
backend and an nginx-served frontend as separate services, with the model
directory mounted read-only and the model cache held in a named volume so that
the aligner is downloaded only once.

\textbf{Single image (production).} A multi-stage Dockerfile builds the frontend
with Node, installs the backend with CPU-only PyTorch wheels --- which avoids
approximately 2\,GB of CUDA payload that would never be used --- optionally
pre-fetches the aligner and Whisper models into the image layer to eliminate a
first-request download, and serves both the API and the compiled static assets
from one Uvicorn process on a single origin. This topology was deployed to a
free-tier container host. Because the API and the UI share an origin, no CORS
configuration is required in production, and the HTTPS origin satisfies the
browser's secure-context requirement for microphone access.

All decision parameters --- error threshold and its bounds, quality-gate limits,
word-verification toggle and threshold, device selection, CORS origins and
upload size limit --- are settable by environment variable, so operating-point
changes require a restart rather than a rebuild (NFR-12).

\section{Testing}
\label{sec:testing}

\subsection{Strategy}

Testing was organised in four levels: unit tests over pure functions,
integration tests over the assembled pipeline, system tests against the running
service and end-to-end tests through the browser. Model quality itself is
evaluated separately in Chapter~\ref{ch:results}; the tests reported here
concern the correctness of the system around the model.

\subsection{Unit testing}

Unit tests cover the deterministic, model-free functions: G2P output and
diacritic stripping; the decode-and-trim path for WAV, FLAC and OGG inputs at
several sample rates; the quality gate against synthetic silence, white noise
and speech; peak-preserving waveform downsampling; the normalisation-constant
resolution order; and threshold clamping.

\subsection{Integration and system testing}

Integration tests exercise the full analysis path on fixed audio, asserting that
the number of returned phonemes matches the G2P output for the target word, that
the spans are monotonically non-decreasing and lie within the trimmed duration,
that probabilities lie in $[0,1]$ and that verdicts are consistent with the
requested threshold. System tests drive the HTTP API directly and assert the
status-code discipline of Section~\ref{sec:api-design}.

The most important integration test is a \emph{parity} test between the notebook
and the service: the same file and target word are processed by both, and the
per-phoneme probabilities are required to agree to within numerical tolerance.
This is the executable form of requirement NFR-03, and it is the test that would
catch a divergence in the normalisation constants or in any feature parameter.

\subsection{Test cases and results}

Table~\ref{tab:testcases} lists the principal test cases and their outcomes.

\begin{table}[H]
\centering
\caption[Test cases and results]{Principal test cases and results.}
\label{tab:testcases}
\scriptsize
\begin{tabularx}{\textwidth}{@{}p{1.15cm}p{2.6cm}Xp{1.5cm}p{1.0cm}@{}}
\toprule
\textbf{ID} & \textbf{Area} & \textbf{Case and expected result} & \textbf{Req.} & \textbf{Result} \\
\midrule
TC-01 & G2P & ``well'' yields \ph{w}\,\ph{ɛ}\,\ph{l}, all in the aligner vocabulary & FR-06 & Pass \\
TC-02 & G2P & Stress and length diacritics absent from the output & FR-06 & Pass \\
TC-03 & Audio & 44.1\,kHz stereo WAV decodes to 16\,kHz mono & FR-05 & Pass \\
TC-04 & Audio & Leading/trailing silence removed; alignment and display share one timebase & FR-07 & Pass \\
TC-05 & Gate & All-zero audio rejected with ``no audio detected'' & FR-14 & Pass \\
TC-06 & Gate & Signal below $-40$\,dBFS rejected as too quiet & FR-14 & Pass \\
TC-07 & Gate & White noise rejected as noise, not analysed & FR-14 & Pass \\
TC-08 & Align & Phoneme count equals G2P output length for the target word & FR-07 & Pass \\
TC-09 & Align & Spans ordered and within the trimmed duration & FR-07 & Pass \\
TC-10 & Model & Probabilities in $[0,1]$; verdict consistent with threshold & FR-08, FR-09 & Pass \\
TC-11 & Model & Notebook and service agree on per-phoneme probabilities & NFR-03 & Pass \\
TC-12 & Verify & Saying ``banana'' for target ``computer'' returns 422 naming what was heard & FR-15 & Pass \\
TC-13 & Verify & Distorted but genuine attempt at the target is not rejected & FR-15 & Pass \\
TC-14 & API & Empty \texttt{target\_word} returns 422 & FR-01 & Pass \\
TC-15 & API & Upload above the size limit returns 413 & NFR-10 & Pass \\
TC-16 & API & Threshold outside $[0.30,0.70]$ is clamped, not rejected & FR-13 & Pass \\
TC-17 & API & \texttt{/api/health} reports model, aligner, Whisper, G2P, device, normalisation & FR-16 & Pass \\
TC-18 & API & With G2P absent, \texttt{/api/analyze} returns 503 and health reports degraded & NFR-09 & Pass \\
TC-19 & Perf & Warm single-word analysis completes within 5\,s on 2 vCPU & NFR-01 & Pass \\
TC-20 & Perf & Models loaded once; second request incurs no load cost & NFR-02 & Pass \\
TC-21 & UI & Record, stop and analyse round trip renders ring, waveform and list & FR-03, FR-12 & Pass \\
TC-22 & UI & Upload path produces the same rendering as the record path & FR-04 & Pass \\
TC-23 & UI & Threshold change re-analyses and updates verdicts & FR-13 & Pass \\
TC-24 & UI & Layout usable at 360\,px viewport width & NFR-07 & Pass \\
TC-25 & UI & Verdicts legible without colour (text label and meter present) & NFR-08 & Pass \\
TC-26 & UI & Backend unreachable shows an ``offline'' indicator, not a silent failure & NFR-13 & Pass \\
\bottomrule
\end{tabularx}
\end{table}

\subsection{Behavioural validation}

Beyond functional testing, one behavioural check was run against the property
the labelling redesign was intended to deliver. A sample of paired recordings
was analysed and the mean fraction of phonemes flagged was computed separately
for the control and the dysarthric member of each pair. Under Method~1 this
comparison had shown no useful separation, with control recordings flagged at
approximately the base rate the labelling scheme had imposed. Under Method~2 the
control recordings returned few or no flagged phonemes while the dysarthric
recordings returned substantially more. A representative single-recording result
is reproduced in Section~\ref{sec:qualitative}.
